\section{Methods}

\subsection{Environment-dependent time-variable galaxy-wide IMF model}

The underlying science framework is an environment-dependent, time-variable galaxy-wide IMF model. IGIMF/GalIMF is used as the computational machinery/tool for the galaxy-wide IMF. In the tests described here, the NGC~1277 relic best-fit parameters are held fixed unless stated otherwise; the scans perturb the enrichment channel rather than refitting the full model.

\subsection{Chemical-evolution setup and standard yields}

The baseline follows the Nomoto-family relic chemical-evolution setup. The analysis tracks surviving stellar mass, light-weighted metallicity proxy, light-weighted \mgfe, $[\mathrm{Si}/\mathrm{Fe}]$, \simg, and internal $M/L$. The current light-weighted abundances are chemical-evolution abundance proxies. They are not full spectral-index abundance inferences.

\subsection{PopIII/PISN-like Si-rich yield-channel tests}

The Si-rich experiments are copied-code yield-channel tests. The mixed PopIII/PISN-like channel is controlled by opt-in environment variables and uses copied yield tables. The representative yield mass used by the current best model is $170\,M_\odot$, but the exact yield table and row remain a manuscript TODO.

\begin{table}
\centering
\caption{Yield-channel parameters for the compact model set. The PopIII/PISN-like mass, fraction, and timing describe copied-code yield-channel tests, not production-code replacements.}
\label{tab:yield_parameters}
\begin{tabular}{lrrrr}
\toprule
Model class & Start & Duration & Mass & Fraction \\
 & Myr & Myr & $M_\odot$ & \\
\midrule
Clean baseline & -- & -- & -- & 0.00 \\
Mild mixed channel & 0 & 30 & 170 & 0.10 \\
Best accepted timed & 0 & 40 & 170 & 0.35 \\
First rejected timed & 0 & 50 & 170 & 0.35 \\
Highest rejected mixed & 0 & 100 & 170 & 0.30 \\
\bottomrule
\end{tabular}
\tablefoot{The exact yield file/table row for the representative $170\,M_\odot$ PopIII/PISN-like channel remains a TODO before manuscript use.}
\end{table}


\subsection{E-MILES NIR-local weighting and index comparison}

E-MILES BaSTI baseFe bimodal SSP spectra were used in two ways. First, NIR-local fluxes around MgI~1.18, SiI~1.20, and SiI~1.59 were used to reweight stellar generations and recompute abundance averages. Second, public baseFe SSP spectra were measured for Eftekhari-style pseudo-equivalent-width indices, including velocity broadening to $\sigma=440\,\mathrm{km\,s^{-1}}$.

These tests do not provide abundance-variable response functions. They test flux weighting and baseline SSP/index behavior, not an independent inference of $[\mathrm{Si}/\mathrm{Fe}]$ or \mgfe.

\subsection{Model acceptance criteria}

\begin{table}
\centering
\caption{Adopted relic-fit acceptance criteria and comparison to the best accepted and first rejected timed models.}
\label{tab:acceptance}
\begin{tabular}{lrrr}
\toprule
Quantity & Tolerance & Best accepted & First rejected \\
\midrule
$|\Delta\mh|$ & $\leq 0.10$ dex & +0.095 & +0.113 \\
$|\Delta\mgfe|$ & $\leq 0.08$ dex & -0.070 & -0.085 \\
$|\Delta\log M|$ & $\leq 0.15$ dex & +0.137 & +0.168 \\
$|\Delta\log M/L|$ & $\leq 0.15$ dex & +0.053 & +0.072 \\
\bottomrule
\end{tabular}
\tablefoot{The best accepted model is \texttt{timed\_s0\_d40\_m170\_f0p35\_e1p00}; the first rejected model is \texttt{timed\_s0\_d50\_m170\_f0p35\_e1p00}.}
\end{table}

